\documentclass[11pt]{article}
\usepackage[margin=1in]{geometry}
\usepackage{amsmath,amssymb}
\usepackage{booktabs}
\usepackage{array}
\usepackage{hyperref}
\hypersetup{hidelinks}

\setlength{\parskip}{4pt plus 1pt minus 1pt}
\setlength{\parindent}{0pt}

\title{Coupling-Channel Structure Predicts Cancer Survival\\
Beyond Mutation Count: Biological Validation\\
on 25{,}776 Patients}
\author{Patrick D.\ McCarthy}
\date{}

\begin{document}
\maketitle

\begin{abstract}
Organizing 509 cancer genes into eight coupling channels---the
signaling systems that maintain multicellular organization---recovers
prognostic signal invisible to flat mutation models. Using only
Cox proportional hazards regression on 25{,}776 metastatic cancer
patients (MSK-MET), we show three results.

(1) Channel count resolves the TMB paradox. Mutation count
becomes \emph{protective} (HR~$= 0.90$, $p < 10^{-32}$) once
channel count is controlled. High mutations with few channels hit
means inefficient cancer. Channel count explains what mutation
count means.

(2) Protein--protein interaction topology independently validates
channel structure. Channels with the lowest PPI coherence
(distributed, shared-service channels) are the most predictive of
survival. Dense wiring provides redundancy; sparse wiring cannot
reroute around breaks.

(3) Channels pair into four symmetric tiers. We test the
complete truth table of logic gates over tier pairs. All four
gates (AND, OR, XOR, NOR) add significant signal---the tiers are
real. The full 2$\times$2 encoding (NOR + AND) improves
prediction in 19 of 22 cancer types (Wilcoxon $p = 0.0001$).

No neural networks. No custom architectures. Standard Cox model,
theory-derived structure, three converging lines of evidence.
The signal is in the structure, not the model.
\end{abstract}

% ------------------------------------------------------------------
\section{Introduction}
\label{sec:intro}
% ------------------------------------------------------------------

Genomic survival prediction in oncology faces a structural
problem. Targeted sequencing panels profile hundreds of genes per
patient, but standard models treat each mutated gene as an
independent binary feature. Cox proportional hazards models
\cite{Cox1972}, random survival forests \cite{Ishwaran2008}, and
deep survival models \cite{Katzman2018} all operate on flat
feature vectors. The question of \emph{which genes to include}
and \emph{how to organize them} is left to feature selection
heuristics or data-driven discovery.

Paper~1 of this series \cite{McCarthy2026a} proves from
information-theoretic first principles that graph structure is
necessary for information preservation under bounded context.
Paper~5 \cite{McCarthy2026e} applies this to cancer:
multicellular organization is maintained by coupling
channels---independent signaling systems that connect cell
behavior to tissue, organ, and organism-level constraints. Cancer
is the predictable consequence of coupling-channel severance.

If coupling channels are the causally relevant structure, then
organizing mutations by channel should recover prognostic signal
that flat models miss. This paper tests that prediction using the
simplest possible model---Cox proportional hazards---so that any
signal detected is attributable to the structure, not the model.

\subsection{Contributions}

\begin{enumerate}
\item A 509-gene, 8-channel mapping organized into 4 hierarchical
   tiers, derived from bounded-context theory and validated against
   PPI topology.
\item Resolution of the TMB paradox: mutation count becomes
   protective once channel count is in the model, decomposing
   the contradictory literature on tumor mutational burden.
\item Independent biological validation of channel structure via
   PPI coherence analysis, showing that distributed channels are
   more predictive than clustered channels.
\item A complete truth-table analysis of tier logic operators
   (AND, OR, XOR, NOR). All four gates add significant signal,
   confirming the tiers are real. The NOR + AND encoding
   (19 of 22 cancer types, $p = 0.0001$) captures the full
   tier state space.
\item Per-cancer-type validation across 22 cancer types,
   identifying the boundary conditions where channel structure
   applies and where it does not.
\end{enumerate}

% ------------------------------------------------------------------
\section{Methods}
\label{sec:methods}
% ------------------------------------------------------------------

\subsection{Coupling-Channel Definitions}

We define eight coupling channels based on the organizational
signaling systems identified in Paper~5 \cite{McCarthy2026e},
organized into four hierarchical tiers. Each channel represents
an independent system by which tissue-level constraints are
communicated to individual cells.

\subsubsection{Tier~0: Cell-Intrinsic}

The core cellular machinery---growth signaling and cell division.

\begin{enumerate}
\item \textbf{PI3K/Growth Signaling:} 174 genes. Mitogenic
   signaling and growth factor response. Key genes include
   \emph{PIK3CA, PTEN, KRAS, NRAS, BRAF, EGFR, ERBB2, MET,
   NF1, NF2, STK11}.
\item \textbf{Cell Cycle:} 95 genes. Growth/quiescence
   checkpoint control. Key genes include \emph{TP53, RB1,
   CDKN2A, CDKN2B, CDK4, CDK6, CCND1, CCNE1, MDM2, MYC}.
\end{enumerate}

\subsubsection{Tier~1: Tissue-Level}

Tissue integrity---DNA repair and structural organization.

\begin{enumerate}
\setcounter{enumi}{2}
\item \textbf{DNA Damage Repair (DDR):} 81 genes. Genomic
   integrity maintenance. Key genes include \emph{BRCA1, BRCA2,
   PALB2, ATM, ATR, CHEK2, RAD51C, MLH1, MSH2, MSH6}.
\item \textbf{Tissue Architecture:} 45 genes. Cell--cell
   adhesion, contact inhibition, and morphogen signaling. Key
   genes include \emph{CDH1, APC, CTNNB1, SMAD4, NOTCH1--4,
   TGFBR1/2}.
\end{enumerate}

\subsubsection{Tier~2: Organism-Level}

Organism-wide coordination---hormonal and immune signaling.

\begin{enumerate}
\setcounter{enumi}{4}
\item \textbf{Endocrine:} 15 genes. Hormone-mediated coupling to
   organism-level physiology. Key genes: \emph{ESR1, PGR, AR,
   FOXA1, GATA3}.
\item \textbf{Immune:} 53 genes. Immune surveillance coupling.
   Key genes include \emph{B2M, HLA-A, HLA-B, JAK1, JAK2,
   STAT1, CD274}.
\end{enumerate}

\subsubsection{Tier~3: Meta-Regulatory}

Epigenetic access control---which genes can be read at all.

\begin{enumerate}
\setcounter{enumi}{6}
\item \textbf{Chromatin Remodeling:} 38 genes. Physical access
   to DNA. Key genes include \emph{ARID1A, ARID1B, SMARCA4,
   SMARCB1, KMT2C, KMT2D}.
\item \textbf{DNA Methylation:} 8 genes. Chemical modification
   state. Key genes include \emph{DNMT3A, DNMT3B, TET2, IDH1,
   IDH2}.
\end{enumerate}

Channel assignments are derived from the signaling biology in
Paper~5 \cite{McCarthy2026e}, cross-referenced with KEGG
\cite{Kanehisa2000}, Reactome \cite{Jassal2020}, and OncoKB
\cite{Chakravarty2017}. The full 509-gene mapping is provided
in the supplementary code.

\subsection{Tier Structure}

The eight channels are organized into four tiers, each containing
a symmetric pair of complementary functions:

\begin{table}[ht]
\centering
\caption{Tier structure: each tier contains two complementary
channels.}
\label{tab:tiers}
\begin{tabular}{@{}clll@{}}
\toprule
Tier & Name & Channel A & Channel B \\
\midrule
0 & Cell-Intrinsic & PI3K/Growth (go) & Cell Cycle (when) \\
1 & Tissue-Level & DDR (fix errors) & Tissue Arch.\ (maintain structure) \\
2 & Organism-Level & Endocrine (receive signals) & Immune (respond to threats) \\
3 & Meta-Regulatory & Chromatin Remodel (read access) & DNA Methylation (write state) \\
\bottomrule
\end{tabular}
\end{table}

These are not arbitrary groupings. Each pair represents
complementary halves of a functional layer: one channel without
the other is incoherent. Growth signaling without cell cycle
control is acceleration without a clock. DNA repair without
tissue architecture is error correction without a filesystem.

This tier structure generates a prediction testable by survival
analysis: tier integrity (both channels intact) should carry
prognostic signal beyond individual channel indicators.

\subsection{Data}

We use the MSK-MET 2021 metastatic cancer cohort
\cite{Nguyen2022}, a publicly available targeted sequencing
dataset with matched survival data:

\begin{itemize}
\item \textbf{Patients:} 25{,}776 metastatic patients. After
   filtering for available overall survival data and at least one
   mutation in the 509-gene channel map: 23{,}371 patients.
\item \textbf{Panel coverage:} 465 of 509 channel genes are
   present on the MSK-MET panel (91.4\%). Per-channel coverage
   ranges from 75\% (DNA Methylation, 6/8 genes) to 94\%
   (PI3K/Growth, DDR).
\item \textbf{Survival endpoint:} Overall survival (OS), with
   time in months and event indicator.
\end{itemize}

\subsection{Feature Engineering}

For each patient, we construct three feature sets of increasing
structural complexity:

\begin{enumerate}
\item \textbf{8 channel indicators} (8 features): binary
   indicator for whether at least one gene in each channel is
   mutated.
\item \textbf{Channel + mutation count} (9 features): channel
   count (number of channels with at least one mutation) plus
   total mutation count.
\item \textbf{Channel + tier state} (16 features): 8 channel
   indicators plus 4 NOR indicators (tier intact: neither channel
   hit) and 4 AND indicators (tier collapsed: both channels hit).
   The full 2$\times$2 tier state per tier.
\end{enumerate}

\subsection{Statistical Methods}

All survival analyses use Cox proportional hazards regression
\cite{Cox1972} with $L_2$ regularization (penalizer $= 0.1$).
Per-cancer-type comparisons use the Wilcoxon signed-rank test
(one-sided) across cancer types with $\geq 50$ patients and
$\geq 10$ events. The sign test and paired $t$-test are reported
as sensitivity analyses.

No neural networks, random forests, or ensemble methods are used.
The deliberate restriction to Cox PH ensures that any detected
signal is attributable to the feature structure, not model
complexity.

% ------------------------------------------------------------------
\section{Results}
\label{sec:results}
% ------------------------------------------------------------------

\subsection{Channel Count Resolves the TMB Paradox}
\label{sec:tmb}

Tumor mutational burden (TMB) has a contradictory literature:
high TMB predicts better immunotherapy response
\cite{Samstein2019} but worse overall survival in some cohorts
and better survival in others. We show that this contradiction
dissolves when mutation count is decomposed by channel.

\begin{table}[ht]
\centering
\caption{Cox regression on MSK-MET ($N = 23{,}371$). Channel
count and mutation count are entered jointly. All models include
$L_2$ regularization.}
\label{tab:tmb}
\begin{tabular}{@{}lccc@{}}
\toprule
Feature & HR & $z$ & $p$ \\
\midrule
Channel count (alone) & 1.07 & --- & $< 10^{-16}$ \\
\midrule
\multicolumn{4}{@{}l}{\emph{Joint model:}} \\
\quad Channel count & 1.09 & $+9.91$ & $\approx 10^{-22}$ \\
\quad Mutation count & 0.90 & $-11.86$ & $\approx 10^{-32}$ \\
\bottomrule
\end{tabular}
\end{table}

Channel count alone is harmful (HR~$= 1.07$): more channels
disrupted means worse prognosis. But once channel count is in the
model, mutation count \emph{flips sign} and becomes protective
(HR~$= 0.90$, $p \approx 10^{-32}$).

The interpretation is direct. At a given number of channels hit,
additional mutations are passengers---evidence of genomic
instability that has not translated into additional channel
severance. High TMB with few channels hit means an inefficient
cancer, mostly accumulating harmless mutations. High TMB with
many channels hit means every mutation landed in a distinct
defense system.

Channel count explains \emph{what mutation count means}. Without
it, mutation count is ambiguous. With it, mutation count measures
cancer inefficiency (inverted).

\subsubsection{Boundary Conditions}

Channel count is significant in 22 of 29 cancer types (75\%).
The 7 types where it fails are cancers driven by mechanisms
outside the channel framework: HPV-driven cancers (anal,
cervical), translocation-driven cancers (sarcomas), environmental
exposure cancers (mesothelioma, hepatobiliary), and single-driver
cancers (GIST). These are ``random'' cancers---damage from
external events---versus the ``tactical'' cancers where defense
systems are sequentially dismantled. The theory predicts its own
domain of applicability.

\subsection{PPI Topology Validates Channel Structure}
\label{sec:ppi}

If channels are biologically real, they should be visible in
protein--protein interaction data independent of our assignment.
We test this using STRING PPI edges (5{,}408 edges among 503
channel genes, confidence $\geq 400$).

For each channel, we compute PPI coherence: the fraction of
possible within-channel gene pairs that have a STRING edge,
compared to the cross-channel baseline. Four distinct channel
types emerge:

\begin{table}[ht]
\centering
\caption{Channel types by PPI coherence and predictive
importance. Importance is the ridge regression coefficient for
per-channel damage features.}
\label{tab:ppi}
\begin{tabular}{@{}llcc@{}}
\toprule
Type & Channels & PPI coherence & Importance \\
\midrule
Physical cluster & DDR (70.6\%), PI3K (79.3\%) & High & Low \\
Execution engine & Cell Cycle (51\%) & Moderate & Near-zero \\
Reactive system & Immune (45\%) & Moderate & Low \\
Shared service & Chromatin (36\%), Tissue (38\%) & Low & \textbf{Highest} \\
Interface layer & Endocrine (16.5\%) & Low & High \\
State layer & DNA Methylation (11.4\%) & Low & High \\
\bottomrule
\end{tabular}
\end{table}

The key finding is paradoxical: channels with the \emph{lowest}
PPI coherence are the \emph{most} predictive of survival.
Tissue Architecture (importance $= 0.0162$) is the single most
important channel damage feature, yet has only 38\% PPI
coherence. Cell Cycle (importance $= 0.0002$) has 51\% coherence
but contributes almost nothing to prediction.

The explanation follows from network redundancy. A PPI cluster
has its repair partners physically nearby: when ATM breaks, ATR
is right there in the same complex. Dense local wiring provides
built-in redundancy. But ARID1A (Chromatin Remodeling) gets
recruited \emph{temporarily} to remodel chromatin at DDR sites,
cell cycle promoters, and immune loci. It is a shared service
called from elsewhere. Break it and there is nothing nearby to
compensate. The recruitment signal goes out and nobody answers.

\begin{quote}
Dense local wiring $=$ redundant $=$ easy to route around $=$
low predictive power.\\
Sparse long-range wiring $=$ no redundancy $=$ impossible to
route around $=$ high predictive power.
\end{quote}

This validates the channel assignments independently of survival
data: the channels are not just statistically useful groupings
but biologically distinct network topologies.

\subsection{Tier Interaction Analysis}
\label{sec:tiers}

The tier structure (Table~\ref{tab:tiers}) predicts that the
\emph{joint state} of both channels in a tier carries information
beyond the individual channel indicators. We test this by
encoding every possible logic gate over each tier's two channel
indicators and comparing per-cancer-type performance.

\subsubsection{Truth Table}

For each tier, the two channel indicators $A$ and $B$ define
four patient states. Each logic gate encodes a different
structural hypothesis:

\begin{table}[ht]
\centering
\caption{Truth table for tier logic gates. $A$ and $B$ are
binary indicators for whether each channel in a tier has at
least one mutation.}
\label{tab:truth}
\begin{tabular}{@{}cc|cccc@{}}
\toprule
$A$ & $B$ & AND & OR & XOR & NOR \\
\midrule
0 & 0 & 0 & 0 & 0 & \textbf{1} \\
0 & 1 & 0 & 1 & 1 & 0 \\
1 & 0 & 0 & 1 & 1 & 0 \\
1 & 1 & \textbf{1} & 1 & 0 & 0 \\
\bottomrule
\end{tabular}
\end{table}

\begin{itemize}
\item \textbf{NOR} (tier intact): neither channel disrupted.
   Hypothesis: an undamaged tier is a definitive protective state.
\item \textbf{AND} (tier collapsed): both channels disrupted.
   Hypothesis: simultaneous loss of complementary functions is
   qualitatively worse than losing either alone.
\item \textbf{XOR} (tier wounded): exactly one channel disrupted.
   Hypothesis: asymmetric load carries a distinct signal.
\item \textbf{OR} (tier touched): at least one channel disrupted.
   The complement of NOR.
\end{itemize}

Note that OR $= 1 - \text{NOR}$, so these two gates encode
identical information. We include both to confirm this
empirically.

\subsubsection{Results by Logic Gate}

We test all four gates. Each model adds 4 tier features (one per
tier) to the 8 channel indicators. All models are fit on the same
22 cancer types where every model converges:

\begin{table}[ht]
\centering
\caption{Per-cancer-type comparison of tier logic gates added to
8 channel indicators ($n = 22$ cancer types where all models
converge). Win/loss/tie counts compare C-index to 8 channels
alone. Wilcoxon $p$ is one-sided (augmented $>$ baseline).}
\label{tab:gates}
\begin{tabular}{@{}lccccc@{}}
\toprule
Model (8ch + 4$\times$gate) & W & L & T & Wilcoxon $p$ & Mean $\Delta$C \\
\midrule
+ AND (collapsed) & 17 & 4 & 1 & 0.0036 & $+0.0045$ \\
+ OR (touched) & 18 & 4 & 0 & 0.0008 & $+0.0049$ \\
+ XOR (wounded) & 16 & 6 & 0 & 0.0008 & $+0.0070$ \\
+ NOR (intact) & 18 & 4 & 0 & 0.0010 & $+0.0049$ \\
\midrule
+ NOR + AND & \textbf{19} & \textbf{3} & \textbf{0} & \textbf{0.0001} & $\mathbf{+0.0077}$ \\
\bottomrule
\end{tabular}
\end{table}

Three observations from Table~\ref{tab:gates}:

\textbf{1. OR and NOR are identical.} Same wins, same losses,
same mean delta, same per-cancer-type results to four decimal
places. This confirms $\text{OR} = 1 - \text{NOR}$---they encode
the same information. We report both to demonstrate this
algebraic fact empirically.

\textbf{2. All four gates add signal.} Every gate achieves
Wilcoxon $p < 0.01$. The tier structure carries prognostic
information regardless of which joint state is encoded. This is
the central finding: the tiers are real.

\textbf{3. NOR + AND dominates.} Combining the two extreme
states---tier intact and tier collapsed---achieves 19--3 at
$p = 0.0001$. This encodes the full 2$\times$2 table: NOR
captures the (0,0) state, AND captures the (1,1) state, and
the wounded states (0,1) and (1,0) are the implicit reference
category, already encoded by the 8 channel indicators.

XOR has the highest mean effect size ($+0.0070$) but the worst
win rate (16--6). It is volatile: large gains in small cancer
types (anal $+0.025$, cervical $+0.034$) offset by more frequent
losses. The wounded tier state is informative but unstable.

\subsubsection{Why NOR Carries the Cleanest Signal}

Among the single-gate models, NOR and OR (which are identical)
achieve the best win rate (18--4) while AND trails (17--4 with
1~tie). The advantage of NOR/OR over AND follows from three
properties:

\textbf{1. Prevalence.} NOR features have higher prevalence than
AND, providing more statistical power per cancer type:

\begin{table}[ht]
\centering
\caption{Feature prevalence: NOR vs.\ AND by tier.}
\label{tab:prevalence}
\begin{tabular}{@{}lcc@{}}
\toprule
Tier & NOR (intact) & AND (collapsed) \\
\midrule
Cell-Intrinsic & 5.0\% & 60.9\% \\
Tissue-Level & 39.3\% & 22.1\% \\
Organism-Level & 73.7\% & 4.8\% \\
Meta-Regulatory & 55.2\% & 4.0\% \\
\bottomrule
\end{tabular}
\end{table}

Tier~0 is the exception: most patients (61\%) have both
cell-intrinsic channels mutated, so AND is common but NOR is
rare. For tiers~2 and~3, AND is so rare ($<5\%$) that per-CT
estimates are unstable, while NOR has sufficient prevalence for
stable coefficients.

\textbf{2. Information content.} NOR encodes a definitive
biological state: the entire tier's defense system is intact.
This is new information that the 8 channel indicators do not
directly encode---knowing that two channels are individually
unmutated does not tell the Cox model that their \emph{joint}
integrity matters. AND, by contrast, is largely redundant with
the two channel indicators both being 1.

\textbf{3. Biological interpretation.} ``This tier is fully
intact'' is a clean prior about system state. ``Both channels in
this tier are broken'' requires both channels to be hit, which
for distributed channels (organism-level, meta-regulatory) is
rare because the cancer would need to waste mutations on
non-core targets. NOR captures the more common and more
informative state: the defense is standing.

\subsubsection{NOR Coefficients}

\begin{table}[ht]
\centering
\caption{NOR coefficients in the 8-channel + 4-NOR Cox model
($N = 23{,}371$).}
\label{tab:nor-coefs}
\begin{tabular}{@{}lccc@{}}
\toprule
Feature & Coefficient & HR & $p$ \\
\midrule
Tissue-level intact & $-0.063$ & 0.94 & 0.019 \\
Organism-level intact & $+0.047$ & 1.05 & 0.16 \\
Cell-intrinsic intact & $-0.003$ & 1.00 & 0.96 \\
Meta-regulatory intact & $-0.028$ & 0.97 & 0.34 \\
\bottomrule
\end{tabular}
\end{table}

Only tissue-level NOR is individually significant ($p = 0.019$,
protective). This is consistent with the PPI analysis: DDR and
Tissue Architecture are the tier whose channels span physical
clusters and shared services. An intact tissue-level tier means
both error correction and structural maintenance are functioning.

The per-cancer-type improvement (18--4 wins) despite weak global
coefficients indicates that NOR's value is context-dependent:
different cancer types load on different tiers, and the NOR
features provide the structural scaffolding for the Cox model to
learn this without explicit cancer-type interaction terms.

% ------------------------------------------------------------------
\section{Discussion}
\label{sec:discussion}
% ------------------------------------------------------------------

\subsection{What Structure Buys You}

The central result is not any single C-index. It is that three
independent lines of evidence---survival decomposition, PPI
topology, and tier interaction analysis---converge on the same
conclusion: coupling channels are biologically real
organizational units, and encoding mutations by channel recovers
prognostic signal that flat models miss.

The TMB paradox resolution is the most striking result. The sign
flip of mutation count from ambiguous to protective ($p \approx
10^{-32}$) once channel count is controlled is not a modeling
trick---it is a decomposition. Total mutations $=$ channel
mutations $+$ passenger mutations. The channels carry the harm;
the passengers carry the noise. Without the decomposition, the
two signals cancel.

\subsection{Distributed Services and Redundancy}

The PPI validation reveals why some channels matter more than
others. Physical clusters (DDR, PI3K/Growth) have built-in
redundancy through dense protein--protein interactions. Shared
services (Chromatin Remodeling, Tissue Architecture) are called
from elsewhere and have no local backup. This maps directly to
engineering: a microservice with no replicas is a single point
of failure.

Cell Cycle (Tier~0) has near-zero predictive importance despite
being the most mutated channel. This is the biological equivalent
of AWS investing maximum redundancy in EC2: the thing everything
depends on is engineered to never fail in a way that matters.
TP53 is the most commonly mutated gene in cancer, yet Cell Cycle
channel damage barely predicts survival---because the cell
invested so heavily in checkpoint redundancy that losing one
component rarely severs the channel.

\subsection{Why Cox PH}

We deliberately restrict all analyses to Cox proportional hazards
regression. This is not a limitation---it is the point. A neural
network that achieves $C = 0.68$ on structured features leaves
open the question of whether the structure or the model
complexity is doing the work. A Cox model that achieves $C =
0.57$ with structure versus $C = 0.50$ without it isolates the
structural contribution.

The per-cancer-type comparisons (Wilcoxon signed-rank tests) are
particularly robust to model choice because they compare the
\emph{same model class} with and without structural features.
The 18--4 NOR result at $p = 0.001$ cannot be attributed to
overfitting, hyperparameter selection, or architecture choices.
It is structure.

\subsection{Boundary Conditions}

Channel structure is significant in 22 of 29 cancer types but
fails in 7. The failures are informative: HPV-driven cancers
(anal, cervical), translocation-driven cancers (sarcomas),
environmental cancers (mesothelioma, hepatobiliary), and
single-driver cancers (GIST). These are cancers where the
damage is not a sequential dismantling of defense
channels---it is a single catastrophic event from outside the
system.

The theory predicts this boundary. Coupling channels are the
structure by which cells maintain organizational constraints.
Cancers that arise from external events (viral integration,
chemical damage, chromosomal rearrangement) bypass the channel
architecture entirely. The 75\% applicability rate is not a
weakness---it is a falsifiable prediction confirmed by the data.

Hepatobiliary cancer is a special case: the liver is the
organism's signal filter. Amplifying edge signals in the liver
means amplifying noise that would normally be cleared. The
channel model underperforms because the organ's function
interferes with the channel abstraction.

\subsection{Limitations}

\textbf{Channel map derived from panel.} The 509-gene channel
map is built from genes on the MSK-IMPACT sequencing panel.
Coverage varies by panel version: 79\% of channel genes are
sequenced on the 2017 panel, 91\% on 2021, and 99\% on the 2026
release. Patients on older panels have missing channel data,
which biases toward null findings.

\textbf{Survival endpoint.} We use overall survival, which is
confounded by treatment, comorbidity, and non-cancer mortality.
Progression-free survival or cancer-specific survival would be
cleaner endpoints but are not uniformly available.

\textbf{No external validation.} All results are on MSK-MET
2021. Replication on TCGA (primary tumors) and MSK-IMPACT 2017
(earlier panel version) is planned but not reported here. The
TCGA replication is particularly important because the TMB
paradox predicts a sign flip between primary and metastatic
cohorts.

\textbf{Tier assignments are theory-driven.} The pairing of
channels into tiers follows from Paper~5's escalation hierarchy,
not from data-driven clustering. This is by design---the point
is to test theoretical predictions---but it means the tier
structure could be improved by empirical optimization. We do not
optimize it, because that would undermine the test.

% ------------------------------------------------------------------
\section{Conclusion}
\label{sec:conclusion}
% ------------------------------------------------------------------

Paper~1 proved that graph structure is necessary for information
preservation. This paper confirms the prediction empirically:
organizing mutations by coupling channel recovers prognostic
signal invisible to flat models.

Three results establish this. Channel count decomposes and
resolves the TMB paradox ($p < 10^{-32}$). PPI topology
independently validates channel assignments as biologically
distinct network structures. The complete truth table of tier
logic gates all add signal---every encoding confirms the tiers
are real---with the full NOR + AND encoding improving 19 of
22~cancer types at $p = 0.0001$ in a standard Cox model with
no tuning.

The signal is in the structure. Not the model, not the
architecture, not the hyperparameters. Five hundred nine genes
organized by coupling-channel theory, fit with the most standard
survival model in biostatistics, recover information that
20{,}000 genes in a flat vector cannot.

Code, data, and the complete 509-gene channel map are available
at [repository URL].

\section*{Acknowledgments}

The author used Claude (Anthropic) for statistical analysis,
literature review, and \LaTeX{} preparation. All theoretical
content, channel definitions, and analytical design are the
author's own.

\bibliographystyle{plain}
\begin{thebibliography}{99}

\bibitem{McCarthy2026a}
P.~D. McCarthy.
\newblock Graph structure is necessary for information preservation under
  bounded context.
\newblock \emph{SSRN preprint}, 2026.

\bibitem{McCarthy2026e}
P.~D. McCarthy.
\newblock Cancer as escalation chain severance under bounded context.
\newblock \emph{SSRN preprint}, 2026.

\bibitem{Cox1972}
D.~R. Cox.
\newblock Regression models and life-tables.
\newblock \emph{Journal of the Royal Statistical Society: Series B},
  34(2):187--220, 1972.

\bibitem{Ishwaran2008}
H.~Ishwaran, U.~B. Kogalur, E.~H. Blackstone, and M.~S. Lauer.
\newblock Random survival forests.
\newblock \emph{The Annals of Applied Statistics}, 2(3):841--860, 2008.

\bibitem{Katzman2018}
J.~L. Katzman, U.~Shaham, A.~Cloninger, J.~Bates, T.~Jiang, and Y.~Kluger.
\newblock {DeepSurv}: personalized treatment recommender system using a {Cox}
  proportional hazards deep neural network.
\newblock \emph{BMC Medical Research Methodology}, 18(1):24, 2018.

\bibitem{Nguyen2022}
B.~Nguyen et~al.
\newblock Genomic characterization of metastatic patterns from prospective
  clinical sequencing of 25,000 patients.
\newblock \emph{Cell}, 185(3):563--575.e11, 2022.

\bibitem{Samstein2019}
R.~M. Samstein, C.-H. Lee, A.~N. Shoushtari, M.~D. Hellmann, R.~Shen,
  Y.~Y. Janjigian, D.~A. Barron, A.~Zehir, E.~J. Jordan, A.~Omuro, and
  others.
\newblock Tumor mutational burden predicts efficacy of {PD-1} blockade across
  cancers.
\newblock \emph{Nature Genetics}, 51(2):202--206, 2019.

\bibitem{Kanehisa2000}
M.~Kanehisa and S.~Goto.
\newblock {KEGG}: {K}yoto encyclopedia of genes and genomes.
\newblock \emph{Nucleic Acids Research}, 28(1):27--30, 2000.

\bibitem{Jassal2020}
B.~Jassal, L.~Matthews, G.~Viteri, C.~Gong, P.~Lorente, A.~Fabregat,
  K.~Sidiropoulos, J.~Cook, M.~Gillespie, R.~Haw, and others.
\newblock The {R}eactome pathway knowledgebase.
\newblock \emph{Nucleic Acids Research}, 48(D1):D498--D503, 2020.

\bibitem{Chakravarty2017}
D.~Chakravarty, J.~Gao, S.~M. Phillips, R.~Kundra, H.~Zhang, J.~Wang,
  J.~E. Rudolph, R.~Yaeger, T.~Soumerai, M.~H. Nissan, and others.
\newblock {OncoKB}: a precision oncology knowledge base.
\newblock \emph{JCO Precision Oncology}, 1:1--16, 2017.

\end{thebibliography}

\end{document}
